\newcommand{\BaslangicTablo}{%
\begin{center}\small
\begin{tabular}{lcccc}
\toprule
fırlatma hücresi & 0 & 6 & 12 & 18 \\
\midrule
başlangıç noktası, eski kural ($\delta=0$): $f_0$ & 1088 & 2729 & $-456$ & 95 \\
başlangıç noktası: $\omega_0$ [rad/s] & $+1.81$ & $-2.48$ & $-0.39$ & $-0.015$ \\
ikinci nokta: $\omega_1$ [rad/s] & $-0.90$ & $+1.2$ & $+0.19$ & $+0.0075$ \\
ara değerleme: $f^*$ & 82.27 & 79.55 & 82.00 & 81.96 \\
\textbf{$\delta^*$'da doğrudan ölçüm: $f$} & \textbf{82.04} & \textbf{81.87} & -- & -- \\
\quad ve orada $\omega_h$ [rad/s] & $-7\times10^{-6}$ & $+2\times10^{-5}$ & -- & -- \\
enerji kuralından başlayarak (Test 1): $f^*$ & 82.04 & 81.90 & 81.89 & 81.94 \\
\bottomrule
\end{tabular}
\end{center}}
\newcommand{\BaslangicYorum}{%
Okunuşu~\BEL:
\begin{enumerate}[nosep]
\item \emph{Donmuş nokta başlangıç kuralından bağımsız bulunuyor.} Eski kuraldan başlanınca bulunan
$\delta^*$ değerleri ($+2.313\times10^{-7}$, $-3.168\times10^{-7}$, $-5.02\times10^{-8}$, $-1.95\times10^{-9}$),
donmuş spin yörüngesinin o hücredeki $|p|$ sapmasıyla (Test~2'nin çapraz kontrolü:
$+2.313\times10^{-7}$, $-3.168\times10^{-7}$, $-5.024\times10^{-8}$, $-1.952\times10^{-9}$) dört basamak aynı.
Yani iki farklı başlangıç kuralı aynı yörüngeyi buluyor.
\item \emph{$\delta^*$'da doğrudan ölçülen $f$, enerji kuralıyla bulunan değerle \%0.04 içinde aynı}
(82.04 ve 81.87'ye karşı 82.04 ve 81.90); orada $\omega_h$ gerçekten sıfır ($10^{-5}$ rad/s).
\item \emph{İki noktalı ara değerleme ise bu testte \%3'e kadar sapıyor} (hücre 6: 79.55). Nedeni: eski
kuralda $\omega_h$'nin $\delta$'ya eğimi farklıdır ($\approx-7.8\times10^{6}$ rad/s; aracın varsaydığı
$-5.24\times10^{6}$ enerji kuralına aittir). Bu yüzden ikinci nokta donmuş noktaya düşmez
($|\omega_1|\approx1$~rad/s) ve ara değerleme geniş bir aralıkta yapılır. Donmuş noktanın dışında $f$'nin
davranışı tam doğrusal olmadığından ara değerleme hata yapar~\HIP.
\item \emph{Önceki bütün üretim koşumları için sonuç~\BEL:} enerji kuralıyla başlayan 666 koşumun hepsinde
ikinci nokta donmuş noktaya çok yakındı ($|\omega_1|$ en fazla $2.7\times10^{-4}$, medyanı $10^{-5}$ rad/s).
Bu testteki 0.9--1.2 rad/s ile karşılaştırınca, önceki sonuçlar fiilen donmuş noktada doğrudan ölçülmüş
sayılır; ara değerleme hatası ihmal edilebilir.
\item \emph{Araçta yapılan iyileştirme:} $|\omega_1|>10^{-3}$ rad/s ise araç artık $\delta^*$'da üçüncü bir ölçüm
yapıyor ve o değeri raporluyor (\kod{W1\_ESIK}). Önceki koşumların hiçbirinde bu eşik aşılmadığı için eski
sonuçlar değişmez; araç başlangıç kuralına karşı sağlam hâle gelir.
\end{enumerate}}

\newcommand{\HamEskiCW}{$-0.08$}
\newcommand{\HamEskiCCW}{$-313$}
\newcommand{\HamEskiC}{150}
\newcommand{\FarkYorum}{%
Okunuşu~\BEL: fark, kapalı yörüngenin büyüklüğüyle birlikte yaklaşık doğrusal büyüyor (kesikli çizgi
eğim 1). Bu, Test~5'teki açıklamayla tutarlı: hata $\propto\sigma\cdot\sigma$, kapalı yörünge
büyüklüğü (rms$_x\cdot$rms$_y$) de $\propto\sigma^2$. Üç rejim:
\begin{itemize}[nosep]
\item \emph{İki demet düzeltmesi} (yörünge $\sim0.01~\mu$m$^2$): fark $\le0.02$, yani sayısal taban; iki
yöntem uyuşuyor.
\item \emph{Tek demet, düzeltilen CCW yönü} ($\sim1~\mu$m$^2$): fark $\le0.17$.
\item \emph{Tek demet CW yönü ve düzeltilmemiş desenler} ($10^2$--$10^4~\mu$m$^2$): fark 1--3000.
Düzeltilmemiş desenlerde eski yöntem $C$'yi 1.1--13 kat yüksek veriyor (medyan 150'ye karşı 49) ve bir
seed'de CW işareti farklı ($-0.08$'e karşı $+38$).
\end{itemize}
Bir incelik: $C$ iki yönün farkıdır. Eski yöntemin hatası iki yönde farklıdır, çünkü iki yönün kapalı
yörüngeleri farklıdır (kuadrupol tekmeleri yönle işaret değiştirir). Bu yüzden hata $C$'de sadeleşmez;
tek demet düzeltmesinde olduğu gibi bir yönün yörüngesi büyük, diğerininki küçükse hata doğrudan $C$'ye
geçer~\YOR.}
